To pass from flowed observables to sharp renormalized local fields, we use the small-flow-time expansion sectorwise
and invert the resulting finite-dimensional mixing system.
Concretely, in each symmetry sector the flowed basis is related to a renormalized local operator basis by a coefficient matrix
$\mathbf C(t,\mu)$.
We require that for sufficiently small $t>0$ this matrix be invertible with a quantitative bound compatible with the SFTE remainder estimates:
\[
\|\mathbf C(t,\mu)^{-1}\|\le C_{\mathrm{inv}}\,\Lambda(t),
\qquad
\Lambda(t)\,t^\eta\to 0\quad (t\downarrow 0).
\]

\paragraph{How to read the widened Lane~A coefficient package.}
There are three logically distinct layers in this section.
\begin{enumerate}[label=(\roman*),leftmargin=2.3em]
\item \emph{Audited sample sectors.}
We begin with explicit scalar, loop, tensor, and concrete dim-$6$ sample blocks.
These are referee-facing calculations showing how the abstract inverse-control mechanism appears in familiar channels.
\item \emph{Canonical exact-dimension blocks.}
For each engineering dimension, CP parity, and Euclidean tensor type, the local block is written in a complete-contraction basis modulo the quotient relations already built into Definition~\ref{def:local_operator_spaces};
the flowed block is obtained by applying the \emph{same contraction formulas} to $G_{\mu\nu}(t,\cdot)$ and flowed covariant derivatives.
Appendix~\ref{app:normalization_crosscheck} then formalizes these blocks as exact-dimension quotient spaces and proves, \emph{before any basis is chosen},
that the small-flow substitution induces the identity on the associated-graded quotient (Theorem~\ref{thm:LaneA_associated_graded_identity}).
The coherent renormalization of Theorem~\ref{thm:field_algebra_truncation} preserves that intrinsic exact-dimension symbol, and only afterwards do we write matrices in the canonical basis.
\item \emph{Finite truncations and full scope.}
Once each exact-dimension diagonal block has quantitative inverse control, the full finite engineering cap follows from lower-triangular assembly,
and Lane~A then reaches the full sharp-local algebra by the inductive-union argument in Section~\ref{sec:sharp_local}.
\end{enumerate}
Appendix~\ref{app:normalization_crosscheck} now records this basis/normalization protocol explicitly.
The explicit scalar/tensor/loop computations therefore remain as normalization cross-checks and audited examples,
while the load-bearing widening step is the canonical finite-truncation theorem near the end of the section.

\medskip
\begin{itemize}
\item CP-even scalar dim-4 (after centering): coefficient $c_S(t,\mu)$ in $\mathcal O_{S,t}\mapsto \mathcal O_S^{\mathrm{ren}}$,
\item CP-odd pseudoscalar dim-4: coefficient $c_Q(t,\mu)$ in $\mathcal O_{Q,t}\mapsto \mathcal O_Q^{\mathrm{ren}}$,
\item small (symmetrized) Wilson loops: coefficient $d_S(t,\mu)$ in $(W_t-1)/\ell^4\mapsto \mathcal O_S^{\mathrm{ren}}$,
\item traceless-symmetric dim-4 tensor probes built from flowed curvature bilinears: finite block $C_T(t,\mu)$ in $\mathbf U_{T,t}\mapsto \mathbf T^{\mathrm{ren}}$,
\item a concrete plaquette/rectangle/chair CP-even dim-$6$ scalar loop block: finite block $D_6^{PRC}(t,\mu)$ in $\mathbf V_{6,t}^{PRC}\mapsto \mathbf O_{6}^{PRC,\mathrm{ren}}$.
\end{itemize}
In each explicitly stabilized case, the inverse is bounded by a \textbf{polylogarithmic} function $\Lambda(t)$,
so $\Lambda(t)t^\eta\to 0$ for every $\eta>0$.
For the concrete dim-$6$ loop block below, we record the finite-dimensional UV/tree-level certificate explicitly:
\[
\begin{aligned}
D_{6,\mathrm{tree}}^{PRC}
&=
\begin{pmatrix}
0 & 1/8\\[0.15em]
1/24 & 0
\end{pmatrix},\\
\det D_{6,\mathrm{tree}}^{PRC}
&=-\frac{1}{192}.
\end{aligned}
\]
and the final theorems in this section then show how the same inverse-control mechanism extends from these audited examples to arbitrary fixed symmetry blocks and hence to every finite sharp-local truncation.

\paragraph{Standing convention (physical Wilson state; chart truncation as an analytic device).}
All expectations are taken in the physical Wilson lattice state.  We occasionally pass through a good-flow chart truncation
to obtain uniform kernel and RG bounds; the truncation error is removed by the de-patching stability lemmas of
Section~\ref{sec:patching} (Lemmas~\ref{lem:depatch_block} and \ref{lem:depatch_continuum}).

In the three target sectors, after the standard subtractions, the mixing matrix is $1\times 1$.

The flowed probe is $\mathcal O_{S,t}=\tfrac14\mathrm{tr}\,G_{\mu\nu}(t)G_{\mu\nu}(t)$.
The only lower-dimension mixing is with the identity; centering removes it:
\[
\mathcal O_{S,t}(f)-\langle \mathcal O_{S,t}(f)\rangle\mathbf 1
= c_S(t,\mu)\,\mathcal O_S^{\mathrm{ren}}(f;\mu) + R_{S,t}^{\circ}(f).
\]

\[
\mathcal O_{Q,t}(f)=c_Q(t,\mu)\,\mathcal O_Q^{\mathrm{ren}}(f;\mu) + R_{Q,t}(f).
\]

For a symmetrized family of flowed plaquettes of side $\ell$ (averaged over orientations) define
\[
\mathcal W_t^{(\ell)}(x):=\frac{1}{6}\sum_{\mu<\nu}\frac{1}{\dim(\rho)}\Re\Tr_{\rho}\, \mathcal P \exp\Big(\oint_{C^{\mu\nu}_\ell} B(t,\cdot)\cdot d x\Big),
\qquad
\ell^2=\kappa_W t.
\]
Then
\[
\frac{\mathcal W_t^{(\ell)}(f)-\langle \mathcal W_t^{(\ell)}(f)\rangle}{\ell^4}
=
d_S(t,\mu)\,\mathcal O_S^{\mathrm{ren}}(f;\mu) + R_{W,t}(f).
\]
Thus in each sector, $\mathbf C(t,\mu)$ reduces to a scalar coefficient and invertibility is simply nonvanishing of that coefficient.

We use the insertion RG machinery from Theorem~\ref{thm:ins_decomp} and Definition~\ref{def:RG_ins}. Fix a renormalization scale $\mu$ in the AF window (so $u(\mu)\le u_\ast$),
and define the ``matching'' step number
\[
n(t,\mu)=\Big\lceil \log_2\Big(\frac{1}{\mu\sqrt t}\Big)\Big\rceil.
\]
The coefficients $c_S(t,\mu)$, $c_Q(t,\mu)$, $d_S(t,\mu)$ are defined as the extracted relevant coefficients at the matching scale
in the insertion RG decomposition ( Theorem~\ref{thm:ins_decomp}).

For each sector, define the one-step multiplier $M_{*,j}$ as follows:
start with the normalized basis insertion $\Phi_{*,j}$ at scale $j$ representing $\mathcal O_*^{\mathrm{ren}}$ in that sector;
apply one good-event (chart-truncated) insertion RG step ( Definition~\ref{def:RG_ins});
extract the coefficient of $\Phi_{*,j+1}$ in the relevant sector at scale $j+1$.
This yields
\[
\Phi_{*,j}\ \xrightarrow{\mathcal R^{\mathrm{ins,good}}_j}\ M_{*,j}\,\Phi_{*,j+1}\ +\ (\text{extracted-irrelevant}).
\]
The number $M_{*,j}$ depends on $(u_j,K_j)$ but is \emph{finite-dimensional}.

\begin{proposition}[One-step multiplier bound]\label{prop:mult_bound}
There exist constants $u_\ast>0$ and $C_{M,S},C_{M,Q},C_{M,W}$ such that whenever the bulk RG lies in the AF ball
($0\le u_j\le u_\ast$ and $\|K_j\|_j\le r_K$),
\begin{equation}\label{eq:M_bounds}
|M_{S,j}-1|\le C_{M,S}\,u_j,\qquad
|M_{Q,j}-1|\le C_{M,Q}\,u_j,\qquad
|M_{W,j}-1|\le C_{M,W}\,u_j.
\end{equation}
We additionally require $u_\ast\le \min\{(2C_{M,S})^{-1},(2C_{M,W})^{-1}\}$ so that $|M_{S,j}|,|M_{W,j}|\ge \tfrac12$.
\end{proposition}

\begin{proof}
A proof (with constants tracked in terms of the insertion integration bound, extraction/projection bounds, and the stable-manifold forcing constant)
is given in Appendix~\ref{app:multiplier_bounds}.
\end{proof}

\begin{remark}[Optional strengthening in the CP-odd channel]
If one inputs additional information about the vanishing of the leading anomalous dimension in the CP-odd channel,
the middle bound in \eqref{eq:M_bounds} can be improved to $|M_{Q,j}-1|=O(u_j^2)$.
The present paper does not use this improvement.
\end{remark}


\paragraph{Comment (status).}
These are finite-dimensional consequences of analyticity of the insertion RG map on the AF ball (as inherited from Theorem~\ref{thm:RG_analytic_bounds}) and normalization $M(0,0)=1$.
Let $u_j$ denote the dyadic running coupling at scale $j$ in the AF window.
\begin{equation}\label{eq:beta}
u_{j+1}=u_j + A\,u_j^2 + r_j,\qquad A:=2b_0\log 2>0,\qquad |r_j|\le B\,u_j^3,
\end{equation}
with $B$ scale-independent.

\begin{lemma}[Sum of couplings grows like $\log(u_{\mathrm{IR}}/u_{\mathrm{UV}})$]\label{lem:sum_u}
Fix $u_\ast>0$ such that $B u_\ast\le A/2$ and assume $u_j\le u_\ast$ for $j=j_0,\dots,j_0+N$.
Then
\begin{equation}\label{eq:sum_u}
\sum_{m=0}^{N-1} u_{j_0+m}
\ \le\ \frac{2}{A_-}\ \log\Big(\frac{u_{j_0+N}}{u_{j_0}}\Big),
\qquad
A_-:=\frac{A}{1+A u_\ast}.
\end{equation}
In particular, if $u_{j_0+N}=u(\mu)$ is fixed and $u_{j_0}=u(\mu_t)\to 0$ as $t\downarrow 0$, then the sum diverges at most like $\log(1/u(\mu_t))$.
\end{lemma}

\begin{proof}
From \eqref{eq:beta} and $B u_\ast\le A/2$, we have $u_{j+1}\ge u_j + (A/2)u_j^2$.
Thus
\[
\frac{1}{u_{j+1}}=\frac{1}{u_j(1+(A/2)u_j)}\le \frac{1}{u_j}-\frac{A/2}{1+(A/2)u_\ast}\le \frac{1}{u_j}-\frac{A_-}{2}.
\]
Iterating yields $u_{j_0+m}\le (u_{j_0}^{-1}-\tfrac{A_-}{2}m)^{-1}$. Summing the harmonic bound gives \eqref{eq:sum_u}.
\end{proof}

Define the RG coefficient at matching scale $\mu$ by
\[
c_S(t,\mu):=\prod_{m=0}^{n(t,\mu)-1} M_{S,j_t+m}\cdot c_{S,\mathrm{UV}}(t),
\]
where $c_{S,\mathrm{UV}}(t)$ is the ``initial'' coefficient at the flow-matching UV scale $j_t$ (by normalization $c_{S,\mathrm{UV}}(t)=1+O(u(\mu_t))$).

\begin{theorem}[CP-even nondegeneracy and polylog inverse]\label{thm:cS}
Assume Proposition~\ref{prop:mult_bound} and that the RG trajectory from scale $\mu_t\simeq 1/\sqrt t$ down to $\mu$ stays in the AF ball ($u_j\le u_\ast$).
Then for all sufficiently small $t$,
\begin{equation}\label{eq:cS_bounds}
0< c_{S,\min}(t,\mu)\ \le\ |c_S(t,\mu)|\ \le\ c_{S,\max}(t,\mu),
\end{equation}
with explicit polylog bounds
\begin{equation}\label{eq:polylog}
|c_S(t,\mu)|^{-1}\ \le\ C_{S,\mathrm{inv}}\ \Big(\frac{u(\mu)}{u(\mu_t)}\Big)^{p_S},
\qquad
p_S:=\frac{4C_{M,S}}{A_-},
\end{equation}
and $C_{S,\mathrm{inv}}$ depending only on the UV normalization bound $|c_{S,\mathrm{UV}}(t)|\ge 1/2$ for small $t$.
In particular, since $u(\mu_t)\downarrow 0$ as $t\downarrow 0$ in an asymptotically free trajectory,
\[
\Lambda_S(t):=|c_S(t,\mu)|^{-1}
\]
grows at most polylogarithmically in $t^{-1}$.
\end{theorem}

\begin{proof}
For $u_j\le u_\ast$ with $u_\ast\le (2C_{M,S})^{-1}$, we have $|M_{S,j}|\ge 1/2$ and
$|\log|M_{S,j}||\le 2|M_{S,j}-1|\le 2C_{M,S}u_j$.
Hence
\[
|\log|c_S(t,\mu)|-\log|c_{S,\mathrm{UV}}(t)||\le 2C_{M,S}\sum_{m=0}^{n-1}u_{j_t+m}.
\]
Use Lemma~\ref{lem:sum_u} and exponentiate, absorbing $|c_{S,\mathrm{UV}}(t)|^{-1}\le 2$ into $C_{S,\mathrm{inv}}$.
\end{proof}

\begin{theorem}[CP-odd nondegeneracy and polylog inverse]\label{thm:cQ}
Assume Proposition~\ref{prop:mult_bound} and that the RG trajectory from scale $\mu_t\simeq 1/\sqrt t$ down to $\mu$ stays in the AF ball ($u_j\le u_\ast$).
Then for all sufficiently small $t$,
\begin{equation}\label{eq:cQ_bounds}
0< c_{Q,\min}(t,\mu)\ \le\ |c_Q(t,\mu)|\ \le\ c_{Q,\max}(t,\mu),
\end{equation}
with explicit polylog bounds
\begin{equation}\label{eq:polylog_Q}
|c_Q(t,\mu)|^{-1}\ \le\ C_{Q,\mathrm{inv}}\ \Big(\frac{u(\mu)}{u(\mu_t)}\Big)^{p_Q},
\qquad
p_Q:=\frac{4C_{M,Q}}{A_-},
\end{equation}
and $C_{Q,\mathrm{inv}}$ depending only on the UV normalization bound $|c_{Q,\mathrm{UV}}(t)|\ge 1/2$ for small $t$.
In particular, $\Lambda_Q(t):=|c_Q(t,\mu)|^{-1}$ grows at most polylogarithmically in $t^{-1}$, hence $\Lambda_Q(t)t^\eta\to0$ for every $\eta>0$.
\end{theorem}

\begin{proof}
From \eqref{eq:M_bounds} we have $|M_{Q,j}-1|\le C_{M,Q}u_j$.
For $u_j\le u_\ast\le (2C_{M,Q})^{-1}$ we have $|\log|M_{Q,j}||\le 2C_{M,Q}u_j$.
Therefore
\[
\Big|\log\Big|\prod_{m=0}^{n-1} M_{Q,j_t+m}\Big|\Big|
\le 2C_{M,Q}\sum_{m=0}^{n-1} u_{j_t+m}.
\]
Apply Lemma~\ref{lem:sum_u} to bound the sum by $\frac{4}{A_-}\log\!\big(\frac{u(\mu)}{u(\mu_t)}\big)$, and absorb the UV normalization factor $c_{Q,\mathrm{UV}}(t)=1+O(u(\mu_t))$ into $C_{Q,\mathrm{inv}}$ for $t$ small.
\end{proof}


At tree level (classical expansion), the symmetrized plaquette probe satisfies
\begin{equation}\label{eq:dS0}
\frac{\mathcal W^{(\ell)}_0(x)-1}{\ell^4}
=\frac{1}{6N}\,\mathcal O_S(x)\ +\ O(\ell^2),
\end{equation}
so the leading coefficient is
\[
d_S^{(0)}=\frac{1}{6N}\neq 0.
\]
Quantum and flow corrections renormalize this by the same RG running as $\mathcal O_S$ plus $O(u)$ perturbations.

\begin{theorem}[Loop coefficient nondegeneracy and inverse bound]\label{thm:dS}
Assume Proposition~\ref{prop:mult_bound} and the AF-ball condition along the trajectory.
Then for sufficiently small $t$,
\[
|d_S(t,\mu)|\ \ge\ \frac{|d_S^{(0)}|}{4}\ \Big(\frac{u(\mu_t)}{u(\mu)}\Big)^{p_W},
\qquad
|d_S(t,\mu)|^{-1}\ \le\ C_{W,\mathrm{inv}}\ \Big(\frac{u(\mu)}{u(\mu_t)}\Big)^{p_W},
\]
with exponent $p_W:=\frac{4C_{M,W}}{A_-}$ and $C_{W,\mathrm{inv}}$ depending only on the UV normalization and the choice of loop family.
In particular, $|d_S(t,\mu)|^{-1}$ grows at most polylogarithmically as $t\downarrow 0$.
\end{theorem}

\begin{proof}
Normalize the UV loop coefficient so that $d_{S,\mathrm{UV}}(t)$ lies within a factor $2$ of $d_S^{(0)}$ for $t$ sufficiently small
(using the classical expansion \eqref{eq:dS0} and $u(\mu_t)\to 0$).
Then run the same multiplier-product bound as in Theorem~\ref{thm:cS} with $M_{W,j}$.
\end{proof}

\begin{corollary}[Sectorwise invertibility condition]\label{cor:Lambda}
In each of the three sectors, set $\Lambda(t)$ to be the corresponding inverse bound:
\[
\Lambda_S(t)=|c_S(t,\mu)|^{-1},\qquad
\Lambda_Q(t)=|c_Q(t,\mu)|^{-1},\qquad
\Lambda_W(t)=|d_S(t,\mu)|^{-1}.
\]
Then $\Lambda(t)$ grows at most polylogarithmically in $t^{-1}$, hence for every $\eta>0$,
\[
\Lambda(t)\,t^\eta\ \xrightarrow[t\downarrow 0]{}\ 0.
\]
Therefore the required small-flow-time invertibility criterion is satisfied sectorwise for the CP-even, CP-odd, and loop generators.
\end{corollary}

\begin{proof}
$t^\eta (\log(1/t))^p\to 0$ as $t\downarrow0$ for all fixed $\eta>0$ and $p\ge 0$.
\end{proof}



\paragraph{A generic finite-dimensional diagonal-block criterion.}
The scalar multiplier-product proofs above use only two ingredients:
near-identity one-step transfer in the AF ball and an invertible UV matching coefficient.
The same mechanism closes any fixed finite symmetry block.
For the exact-dimension quotient blocks attached to Definition~\ref{def:local_operator_spaces}, the needed coefficient extraction is now theorem-level rather than schematic:
Appendix~\ref{app:normalization_crosscheck} proves that quotienting by lower-dimension terms and homogeneous relations commutes with the small-flow substitution and yields a literal identity matrix on every arbitrary exact-dimension block.
Thus the proposition below is applied to canonically defined finite-dimensional blocks, not to a sector-by-sector template.

\begin{proposition}[Finite-dimensional one-step block multiplier bound]\label{prop:block_mult_bound}
Let $\mathfrak s$ be a fixed finite symmetry sector with basis size $m_{\mathfrak s}<\infty$
(for instance a traceless-symmetric tensor channel or the concrete plaquette/rectangle/chair dim-$6$ scalar block treated below).
Let
\[
\mathbf\Phi_{\mathfrak s,j}
=
\bigl(\Phi_{\mathfrak s,j}^{(1)},\dots,\Phi_{\mathfrak s,j}^{(m_{\mathfrak s})}\bigr)^{\mathsf T}
\]
be the normalized basis insertion vector at scale $j$ and define the extracted one-step transfer matrix
$M_{\mathfrak s,j}\in \mathrm{Mat}_{m_{\mathfrak s}}(\mathbb R)$ by
\[
\mathbf\Phi_{\mathfrak s,j}
\xrightarrow{\mathcal R^{\mathrm{ins,good}}_j}
M_{\mathfrak s,j}\,\mathbf\Phi_{\mathfrak s,j+1}
+ (\text{extracted-irrelevant}).
\]
Then there exists a constant $C_{M,\mathfrak s}<\infty$ such that whenever the bulk RG lies in the AF ball
($0\le u_j\le u_\ast$ and $\|K_j\|_j\le r_K$),
\begin{equation}\label{eq:block_mult_near_identity}
\|M_{\mathfrak s,j}-I\|_{\mathrm{op}}\le C_{M,\mathfrak s}\,u_j.
\end{equation}
After shrinking $u_\ast$ so that $C_{M,\mathfrak s}u_\ast\le \tfrac12$, each $M_{\mathfrak s,j}$ is invertible and
\begin{equation}\label{eq:block_mult_inverse_bound}
\|M_{\mathfrak s,j}^{-1}\|_{\mathrm{op}}
\le (1-C_{M,\mathfrak s}u_j)^{-1}
\le 1+2C_{M,\mathfrak s}u_j
\le \exp(2C_{M,\mathfrak s}u_j).
\end{equation}
\end{proposition}

\begin{proof}
For each input basis vector $\Phi_{\mathfrak s,j}^{(q)}$ and each output coefficient index $p$,
let $\mathsf P^{(p)}_{\mathfrak s,j+1}$ denote the bounded projection extracting the coefficient of
$\Phi_{\mathfrak s,j+1}^{(p)}$ after one insertion RG step.
Exactly as in Appendix~\ref{app:multiplier_bounds}, Lemma~\ref{lem:ins_step_Lip} gives
\[
\big\|\mathcal R^{\mathrm{ins,good}}_j(u,K,\Phi_{\mathfrak s,j}^{(q)})
-\mathcal R^{\mathrm{ins,good}}_j(0,0,\Phi_{\mathfrak s,j}^{(q)})\big\|_{j+1}
\le C_{\mathrm{Lip},\mathfrak s}\,\bigl(u+\|K\|_j\bigr),
\]
with a constant uniform in $j$ because the sector dimension is fixed.
Hence
\[
\big|(M_{\mathfrak s,j}-I)_{pq}\big|
\le C_{\mathrm{proj},\mathfrak s}\,C_{\mathrm{Lip},\mathfrak s}\,\bigl(u_j+\|K_j\|_j\bigr).
\]
On the AF ball the stable-manifold forcing bound gives
$\|K_j\|_j\le \frac{C_F}{1-\kappa}\,u_\ast\,u_j$.
Because $m_{\mathfrak s}<\infty$, operator norm and max-entry norm are equivalent, so after absorbing the fixed sector-size factor we obtain
\eqref{eq:block_mult_near_identity}.
If $\|M_{\mathfrak s,j}-I\|_{\mathrm{op}}\le \tfrac12$, the Neumann-series bound gives
$\|M_{\mathfrak s,j}^{-1}\|_{\mathrm{op}}\le (1-\|M_{\mathfrak s,j}-I\|_{\mathrm{op}})^{-1}$,
which implies \eqref{eq:block_mult_inverse_bound}.
\end{proof}

\begin{theorem}[Generic finite-dimensional diagonal-block nondegeneracy]\label{thm:block_polylog}
Let $\mathfrak s$ be a fixed finite symmetry block and let
\[
C_{\mathfrak s}(t,\mu)\in \mathrm{Mat}_{m_{\mathfrak s}}(\mathbb R)
\]
be the coefficient matrix in the small-flow-time expansion from flowed probes to renormalized local operators in that block.
Assume:
\begin{enumerate}[label=(\roman*),leftmargin=2.3em]
\item the RG trajectory from the matching scale $\mu_t\simeq 1/\sqrt t$ down to $\mu$ stays in the AF ball;
\item the one-step transfer matrices satisfy Proposition~\ref{prop:block_mult_bound};
\item the UV matching matrix has a nondegenerate limit,
\begin{equation}\label{eq:block_UV_nondeg}
C_{\mathfrak s,\mathrm{UV}}(t)=C_{\mathfrak s}^{(0)}+O\!\big(u(\mu_t)\big),
\qquad
C_{\mathfrak s}^{(0)}\in \mathrm{GL}_{m_{\mathfrak s}}(\mathbb R).
\end{equation}
\end{enumerate}
Then for all sufficiently small $t$ the block $C_{\mathfrak s}(t,\mu)$ is invertible and obeys the polylogarithmic inverse bound
\begin{equation}\label{eq:block_polylog_bound}
\|C_{\mathfrak s}(t,\mu)^{-1}\|_{\mathrm{op}}
\le
C_{\mathfrak s,\mathrm{inv}}
\Big(\frac{u(\mu)}{u(\mu_t)}\Big)^{p_{\mathfrak s}},
\qquad
p_{\mathfrak s}:=\frac{4C_{M,\mathfrak s}}{A_-},
\end{equation}
for a constant $C_{\mathfrak s,\mathrm{inv}}<\infty$ depending only on
$\|(C_{\mathfrak s}^{(0)})^{-1}\|_{\mathrm{op}}$ and the small-$t$ UV normalization.
In particular, $\|C_{\mathfrak s}(t,\mu)^{-1}\|_{\mathrm{op}}$ grows at most polylogarithmically as $t\downarrow 0$.
\end{theorem}

\begin{proof}
Write $n=n(t,\mu)$ and factor the coefficient block as
\[
C_{\mathfrak s}(t,\mu)
=
M_{\mathfrak s,j_t+n-1}\cdots M_{\mathfrak s,j_t}\,C_{\mathfrak s,\mathrm{UV}}(t).
\]
By \eqref{eq:block_UV_nondeg}, $C_{\mathfrak s,\mathrm{UV}}(t)$ is invertible for all sufficiently small $t$ and
$\|C_{\mathfrak s,\mathrm{UV}}(t)^{-1}\|_{\mathrm{op}}\le 2\|(C_{\mathfrak s}^{(0)})^{-1}\|_{\mathrm{op}}$ after shrinking $t_0$ if necessary.
Therefore
\[
\|C_{\mathfrak s}(t,\mu)^{-1}\|_{\mathrm{op}}
\le
\|C_{\mathfrak s,\mathrm{UV}}(t)^{-1}\|_{\mathrm{op}}
\prod_{m=0}^{n-1}\|M_{\mathfrak s,j_t+m}^{-1}\|_{\mathrm{op}}.
\]
Using \eqref{eq:block_mult_inverse_bound},
\[
\prod_{m=0}^{n-1}\|M_{\mathfrak s,j_t+m}^{-1}\|_{\mathrm{op}}
\le
\exp\!\Big(2C_{M,\mathfrak s}\sum_{m=0}^{n-1}u_{j_t+m}\Big).
\]
Apply Lemma~\ref{lem:sum_u} to bound the sum by
$\frac{2}{A_-}\log\!\big(\frac{u(\mu)}{u(\mu_t)}\big)$ and absorb the UV prefactor into
$C_{\mathfrak s,\mathrm{inv}}$.
This yields \eqref{eq:block_polylog_bound}.
\end{proof}

\paragraph{New diagonal block: traceless-symmetric dim-4 tensor channel.}
Define the flowed traceless tensor probe
\[
U_{\mu\nu,t}(x)
:=
\tr\!\Big(
G_{\mu\lambda}(t,x)G_{\nu\lambda}(t,x)
-\frac14\delta_{\mu\nu}\,G_{\rho\sigma}(t,x)G_{\rho\sigma}(t,x)
\Big).
\]
Smearing against traceless symmetric test tensors $h^{\mu\nu}$ gives a finite family $\mathbf U_{T,t}(h)$.
Because the channel is traceless symmetric, there is no mixing with the identity or the scalar $F^2$ sector;
the diagonal block is therefore the renormalized traceless-tensor family $\mathbf T^{\mathrm{ren}}(h;\mu)$ itself:
\[
\mathbf U_{T,t}(h)
=
C_T(t,\mu)\,\mathbf T^{\mathrm{ren}}(h;\mu)+\mathbf R_{T,t}(h).
\]
By normalization of the tensor basis at the Gaussian point, the UV matching matrix satisfies
\[
C_{T,\mathrm{UV}}(t)=I_{m_T}+O\!\big(u(\mu_t)\big).
\]

\begin{theorem}[Traceless-tensor diagonal-block inverse control]\label{thm:cT_block}
Assume the AF-ball hypothesis along the trajectory and the one-step block multiplier bound of
Proposition~\ref{prop:block_mult_bound} in the traceless-symmetric dim-4 tensor channel.
Then for all sufficiently small $t$,
\begin{equation}\label{eq:cT_polylog}
\|C_T(t,\mu)^{-1}\|_{\mathrm{op}}
\le
C_{T,\mathrm{inv}}
\Big(\frac{u(\mu)}{u(\mu_t)}\Big)^{p_T},
\qquad
p_T:=\frac{4C_{M,T}}{A_-},
\end{equation}
so $\Lambda_T(t):=\|C_T(t,\mu)^{-1}\|_{\mathrm{op}}$ grows at most polylogarithmically in $t^{-1}$.
In particular, $\Lambda_T(t)\,t^\eta\to 0$ for every $\eta>0$.
\end{theorem}

\begin{proof}
Apply Theorem~\ref{thm:block_polylog} with $\mathfrak s=T$ and $C_T^{(0)}=I_{m_T}$.
\end{proof}

\paragraph{New diagonal block: concrete plaquette/rectangle/chair dim-$6$ scalar loop sector.}
Let $X_{P,t}(f)$, $X_{R,t}(f)$, and $X_{C,t}(f)$ denote centered, orientation-averaged flowed loop densities supported in a fixed $2\ell\times 2\ell$ box, with
\[
\ell^2=\kappa_W t,
\]
corresponding respectively to the plaquette, the planar $1\times 2$ rectangle, and the chair loop.
Normalize the three families so that their common tree-level dim-$4$ scalar coefficient is the same constant $d_S^{(0)}$.
Define the improved CP-even dim-$6$ probes
\[
L_{6,+,t}^{(1)}(f):=\ell^{-2}\bigl(X_{R,t}(f)-X_{P,t}(f)\bigr),
\qquad
L_{6,+,t}^{(2)}(f):=\ell^{-2}\bigl(X_{C,t}(f)-X_{P,t}(f)\bigr).
\]
Because the dim-$4$ contribution cancels by construction, both probes have leading engineering dimension $6$ and live in the CP-even scalar sector.
Write
\[
\mathbf V_{6,t}^{PRC}(f):=
\begin{pmatrix}
L_{6,+,t}^{(1)}(f)\\[0.35em]
L_{6,+,t}^{(2)}(f)
\end{pmatrix}.
\]


\begin{proposition}[Finite UV determinant certificate for the plaquette/rectangle/chair dim-$6$ block]\label{prop:D6_tree_certificate}
For this audited sample block choose the on-shell CP-even dim-$6$ gauge basis
\[
\mathbf O_{6}^{\mathrm{ren}}(f;\mu)
:=
\begin{pmatrix}
\mathcal G_{6}^{(1),\mathrm{ren}}(f;\mu)\\[0.35em]
\mathcal G_{6}^{(2),\mathrm{ren}}(f;\mu)
\end{pmatrix},
\]
where
\[
\mathcal G_{6}^{(1)}\sim \sum_{\mu,\nu,\rho}\Tr\big([D_\mu,F_{\nu\rho}]\,[D_\mu,F_{\nu\rho}]\big),
\qquad
\mathcal G_{6}^{(2)}\sim \sum_{\mu,\nu}\Tr\big([D_\mu,F_{\mu\nu}]\,[D_\mu,F_{\mu\nu}]\big).
\]
This on-shell gauge basis is used only for the explicit certificate; Corollary~\ref{cor:D6_basis_transfer} later transports the same inverse control to any other fixed basis in the same dim-$6$ scalar block.
Then the tree-level matching of the improved probes against that basis has the form
\begin{equation}\label{eq:D6_tree_matrix}
\mathbf V_{6,t}^{PRC}(f)
=
D_{6,\mathrm{tree}}^{PRC}\,\mathbf O_{6}^{\mathrm{ren}}(f;\mu_t)
+\mathbf r_{6,t}^{PRC}(f),
\qquad
\mu_t\simeq 1/\sqrt t,
\end{equation}
with
\[
\|\mathbf r_{6,t}^{PRC}(f)\|=O(\ell^2)+O\!\big(u(\mu_t)\big)
\]
and with the explicit fixed matrix
\begin{equation}\label{eq:D6_tree_det_matrix}
D_{6,\mathrm{tree}}^{PRC}
=
\begin{pmatrix}
0 & \frac18\\[0.45em]
\frac1{24} & 0
\end{pmatrix}.
\end{equation}
Hence the finite UV determinant certificate is
\begin{equation}\label{eq:D6_tree_det}
\Delta_{PRC}:=\det D_{6,\mathrm{tree}}^{PRC}
=
-\frac{1}{192}\neq 0.
\end{equation}
In particular, $D_{6,\mathrm{tree}}^{PRC}\in GL_2(\mathbb R)$.
\end{proposition}

\begin{proof}
Use the standard tree-level Symanzik / L\"uscher--Weisz local matching formulas for the normalized four-link / six-link gauge-loop basis \cite{LuscherWeisz1985,Husung2023}:
for loop coefficients $(e_0,e_1,e_2,e_3)$ satisfying
\[
e_0+8e_1+8e_2+16e_3=1,
\]
the two CP-even dim-$6$ gauge coefficients in the on-shell basis above are
\[
\bar c_1=\frac{e_2}{3},
\qquad
\bar c_2=e_1-e_3+\frac1{12}.
\]
After the common dim-$4$ normalization adopted for $X_{P,t},X_{R,t},X_{C,t}$, the plaquette, rectangle, and chair densities correspond respectively to
\[
P:(e_0,e_1,e_2,e_3)=(1,0,0,0),
\qquad
R:(0,\tfrac18,0,0),
\qquad
C:(0,0,\tfrac18,0).
\]
Hence their tree-level dim-$6$ coefficients are
\[
P:(\bar c_1,\bar c_2)=\Bigl(0,\frac1{12}\Bigr),
\qquad
R:(\bar c_1,\bar c_2)=\Bigl(0,\frac5{24}\Bigr),
\qquad
C:(\bar c_1,\bar c_2)=\Bigl(\frac1{24},\frac1{12}\Bigr).
\]
At matching scale $\mu_t\simeq t^{-1/2}$ this gives
\[
X_{P,t}(f)=d_S^{(0)}\,\mathcal O_S^{\mathrm{ren}}(f;\mu_t)+\ell^2\Bigl(\frac1{12}\,\mathcal G_{6}^{(2),\mathrm{ren}}(f;\mu_t)\Bigr)+r_{P,t}(f),
\]
\[
X_{R,t}(f)=d_S^{(0)}\,\mathcal O_S^{\mathrm{ren}}(f;\mu_t)+\ell^2\Bigl(\frac5{24}\,\mathcal G_{6}^{(2),\mathrm{ren}}(f;\mu_t)\Bigr)+r_{R,t}(f),
\]
\[
X_{C,t}(f)=d_S^{(0)}\,\mathcal O_S^{\mathrm{ren}}(f;\mu_t)+\ell^2\Bigl(\frac1{24}\,\mathcal G_{6}^{(1),\mathrm{ren}}(f;\mu_t)+\frac1{12}\,\mathcal G_{6}^{(2),\mathrm{ren}}(f;\mu_t)\Bigr)+r_{C,t}(f),
\]
with each remainder $r_{*,t}(f)=O(\ell^4)+O(\ell^2u(\mu_t))$.
Subtracting the common plaquette term and dividing by $\ell^2$ yields
\[
L_{6,+,t}^{(1)}(f)=\frac18\,\mathcal G_{6}^{(2),\mathrm{ren}}(f;\mu_t)+O(\ell^2)+O\!\big(u(\mu_t)\big),
\]
\[
L_{6,+,t}^{(2)}(f)=\frac1{24}\,\mathcal G_{6}^{(1),\mathrm{ren}}(f;\mu_t)+O(\ell^2)+O\!\big(u(\mu_t)\big).
\]
This is exactly \eqref{eq:D6_tree_matrix} with matrix \eqref{eq:D6_tree_det_matrix}; taking the determinant gives \eqref{eq:D6_tree_det}.
\end{proof}

Choose the loop-adapted renormalized basis
\[
\mathbf O_{6}^{PRC,\mathrm{ren}}(f;\mu)
:=
D_{6,\mathrm{tree}}^{PRC}\,\mathbf O_{6}^{\mathrm{ren}}(f;\mu).
\]
Because $D_{6,\mathrm{tree}}^{PRC}$ is fixed and invertible, this is just a $t$-independent finite change of basis.
In that basis the UV matching block is normalized to the identity at tree level:
\[
\mathbf V_{6,t}^{PRC}(f)
=
D_6^{PRC}(t,\mu)\,\mathbf O_{6}^{PRC,\mathrm{ren}}(f;\mu)+\mathbf R_{6,t}^{PRC}(f),
\qquad
D_{6,\mathrm{UV}}^{PRC}(t)=I_2+E_{6}^{PRC}(t),
\]
with
\[
\|E_{6}^{PRC}(t)\|_{\mathrm{op}}\le C_{6,\mathrm{UV}}^{PRC}\,u(\mu_t)
\]
for all sufficiently small $t$.

\begin{theorem}[Concrete plaquette/rectangle/chair dim-$6$ block inverse control]\label{thm:D6_block}
Assume:
\begin{enumerate}[label=(\roman*),leftmargin=2.3em]
\item the AF-ball hypothesis holds along the matching trajectory;
\item the one-step block multiplier bound of Proposition~\ref{prop:block_mult_bound} holds in this dim-$6$ scalar block.
\end{enumerate}
Then, after shrinking $t_0$ so that $C_{6,\mathrm{UV}}^{PRC}u(\mu_t)\le \tfrac12$ on $(0,t_0]$, the UV block is uniformly invertible with
\begin{equation}\label{eq:D6_UV_inv}
\|(D_{6,\mathrm{UV}}^{PRC}(t))^{-1}\|_{\mathrm{op}}\le 2.
\end{equation}
Consequently, for all sufficiently small $t$,
\begin{equation}\label{eq:D6_polylog}
\|D_6^{PRC}(t,\mu)^{-1}\|_{\mathrm{op}}
\le
C_{6,\mathrm{inv}}
\Big(\frac{u(\mu)}{u(\mu_t)}\Big)^{p_6},
\qquad
p_6:=\frac{4C_{M,6}}{A_-},
\end{equation}
so the inverse control in this block is again polylogarithmic.
In particular, $\Lambda_6(t):=\|D_6^{PRC}(t,\mu)^{-1}\|_{\mathrm{op}}$ satisfies $\Lambda_6(t)t^\eta\to 0$ for every $\eta>0$.
\end{theorem}

\begin{proof}
The uniform UV inverse bound \eqref{eq:D6_UV_inv} is the Neumann-series estimate for $I_2+E_{6}^{PRC}(t)$ once $\|E_{6}^{PRC}(t)\|_{\mathrm{op}}\le \tfrac12$.
Now apply Theorem~\ref{thm:block_polylog} with $\mathfrak s=6$ and with UV normalization constant coming from \eqref{eq:D6_UV_inv}.
\end{proof}

\begin{corollary}[Basis transfer for the explicit dim-$6$ loop block]\label{cor:D6_basis_transfer}
Let $B\in GL_2(\mathbb R)$ be any fixed change of basis on the same CP-even dim-$6$ scalar sector and set
\[
D_6^{B}(t,\mu):=B^{-1}D_6^{PRC}(t,\mu)B.
\]
Then
\[
\|D_6^{B}(t,\mu)^{-1}\|_{\mathrm{op}}
\le
\|B^{-1}\|_{\mathrm{op}}\|B\|_{\mathrm{op}}\,C_{6,\mathrm{inv}}
\Big(\frac{u(\mu)}{u(\mu_t)}\Big)^{p_6},
\]
so every fixed basis transported from the concrete plaquette/rectangle/chair block inherits the same polylogarithmic inverse control.
\end{corollary}

\begin{proof}
This is immediate from
\[
(D_6^{B}(t,\mu))^{-1}=B^{-1}(D_6^{PRC}(t,\mu))^{-1}B
\]
and Theorem~\ref{thm:D6_block}.
\end{proof}

\begin{corollary}[Diagonal-block inverse control beyond the scalar $1\times 1$ list]\label{cor:Lambda_extended}
The required small-flow-time invertibility criterion is now available in the following additional diagonal blocks:
\begin{enumerate}[label=(\alph*),leftmargin=2.3em]
\item the traceless-symmetric dim-4 tensor block, with inverse control $\Lambda_T(t)=\|C_T(t,\mu)^{-1}\|_{\mathrm{op}}$ from
Theorem~\ref{thm:cT_block};
\item the concrete plaquette/rectangle/chair CP-even dim-$6$ scalar loop block, with inverse control
\[
\Lambda_6(t)=\|D_6^{PRC}(t,\mu)^{-1}\|_{\mathrm{op}}
\]
from Theorem~\ref{thm:D6_block}, and likewise any fixed basis transported from it by Corollary~\ref{cor:D6_basis_transfer};
\item any finite direct sum of the scalar $1\times 1$ sectors and the new finite blocks, after applying the lower-triangular assembly result below.
\end{enumerate}
In each case the inverse growth is at most polylogarithmic in $t^{-1}$.
\end{corollary}

\begin{proof}
Parts (a) and (b) are exactly Theorems~\ref{thm:cT_block} and \ref{thm:D6_block}, together with Corollary~\ref{cor:D6_basis_transfer}.
Part (c) follows by combining those theorems with Proposition~\ref{prop:LaneA_triangular_assembly} below.
\end{proof}
\paragraph{From diagonal blocks to finite lower-triangular families.}
The graded triangular structure expected from power counting is already built into the abstract local-field construction
(Theorem~\ref{thm:field_algebra_truncation}).
For Lane~A one therefore needs not only isolated diagonal-block nondegeneracy, but also a finite-dimensional assembly statement
showing that once the diagonal blocks are controlled, the whole finite lower-triangular system is controlled as well.

\begin{proposition}[Finite lower-triangular assembly of inverse control]\label{prop:LaneA_triangular_assembly}
Let a finite family of flowed probes / renormalized local operators be ordered so that the corresponding coefficient matrix
$\mathbf C(t,\mu)$ is block lower triangular:
\[
\mathbf C(t,\mu)=\bigl(C_{ab}(t,\mu)\bigr)_{1\le a,b\le r},
\qquad
C_{ab}(t,\mu)=0 \ \text{ for } a<b,
\]
with $r<\infty$.
Write
\[
\mathbf D(t,\mu):=\mathrm{diag}\bigl(C_{11}(t,\mu),\dots,C_{rr}(t,\mu)\bigr),
\qquad
\mathbf L(t,\mu):=\mathbf C(t,\mu)-\mathbf D(t,\mu).
\]
Assume:
\begin{enumerate}[label=(\roman*),leftmargin=2.3em]
\item each diagonal block $C_{aa}(t,\mu)$ is invertible and obeys
\[
\|C_{aa}(t,\mu)^{-1}\|\le \Lambda_a(t,\mu);
\]
\item each off-diagonal block is uniformly bounded for $0<t\le t_0$,
\[
\|C_{ab}(t,\mu)\|\le M_{ab}\qquad (b<a);
\]
\item the number of blocks and their sizes are fixed.
\end{enumerate}
Set
\[
\Lambda_{\max}(t,\mu):=\max_{1\le a\le r}\Lambda_a(t,\mu),
\qquad
M_{\mathrm{off}}:=\max_{b<a} M_{ab}.
\]
Then $\mathbf C(t,\mu)$ is invertible for $0<t\le t_0$ and there exists a constant $C_{\mathrm{tri}}<\infty$,
depending only on the block sizes and the finitely many $M_{ab}$, such that
\begin{equation}\label{eq:LaneA_triangular_inverse_bound}
\|\mathbf C(t,\mu)^{-1}\|
\le
C_{\mathrm{tri}}\,(1+\Lambda_{\max}(t,\mu))^{r-1}\,\Lambda_{\max}(t,\mu).
\end{equation}
By finite-dimensional norm equivalence, the same bound holds (after adjusting $C_{\mathrm{tri}}$) for the max-entry norm
$\max_{i,j}|(\mathbf C(t,\mu)^{-1})_{ij}|$ used in Theorem~\ref{thm:LaneA_closure}.
\end{proposition}

\begin{proof}
Factor
\[
\mathbf C(t,\mu)=\mathbf D(t,\mu)\,\bigl(I+\mathbf N(t,\mu)\bigr),
\qquad
\mathbf N(t,\mu):=\mathbf D(t,\mu)^{-1}\mathbf L(t,\mu).
\]
Because $\mathbf N(t,\mu)$ is strictly block lower triangular with $r$ block rows, it is nilpotent of order at most $r$:
\[
\mathbf N(t,\mu)^r=0.
\]
Hence
\[
\bigl(I+\mathbf N(t,\mu)\bigr)^{-1}
=
\sum_{k=0}^{r-1}(-\mathbf N(t,\mu))^k.
\]
Moreover,
\[
\|\mathbf D(t,\mu)^{-1}\|
\le C_{\mathrm{diag}}\,\Lambda_{\max}(t,\mu)
\]
for a constant $C_{\mathrm{diag}}$ depending only on the fixed block sizes and the chosen matrix norm.
Since each off-diagonal block of $\mathbf L$ is uniformly bounded and the block sizes are fixed,
there exists $C_{\mathrm{off}}<\infty$ such that
\[
\|\mathbf N(t,\mu)\|\le C_{\mathrm{off}}\,\Lambda_{\max}(t,\mu).
\]
Therefore
\[
\|\mathbf C(t,\mu)^{-1}\|
\le
\|\mathbf D(t,\mu)^{-1}\|\,
\sum_{k=0}^{r-1}\|\mathbf N(t,\mu)\|^k
\le
C_{\mathrm{diag}}\,\Lambda_{\max}(t,\mu)\sum_{k=0}^{r-1}\bigl(C_{\mathrm{off}}\Lambda_{\max}(t,\mu)\bigr)^k,
\]
which is bounded by the right-hand side of \eqref{eq:LaneA_triangular_inverse_bound} after absorbing the finite geometric sum into
$C_{\mathrm{tri}}(1+\Lambda_{\max})^{r-1}$.
\end{proof}

\begin{corollary}[Finite mixed-sector closure from diagonal-block control]\label{cor:LaneA_mixed_sector}
Consider any finite family of renormalized local operators whose small-flow-time coefficient matrix is block lower triangular
(as predicted by power counting / Theorem~\ref{thm:field_algebra_truncation}).
If each diagonal block admits a quantitative inverse bound $\Lambda_a(t,\mu)$ and each off-diagonal block is uniformly bounded,
then the full family satisfies the Lane~A inverse-control hypothesis \eqref{eq:LaneA_quant_inverse_bound} with
\[
\Lambda_{\mathrm{inv}}^{\mathrm{full}}(t,\mu)
:=
C_{\mathrm{tri}}\,(1+\Lambda_{\max}(t,\mu))^{r-1}\,\Lambda_{\max}(t,\mu).
\]
In particular, whenever the diagonal-block controls are polylogarithmic in $t^{-1}$, the assembled control
$\Lambda_{\mathrm{inv}}^{\mathrm{full}}(t,\mu)$ is again polylogarithmic.
Thus Lane~A applies not only sectorwise, but also to any finite direct sum / lower-triangular assembly built from the presently stabilized
CP-even scalar, CP-odd scalar, small-loop, and traceless-tensor diagonal blocks, together with the concrete plaquette/rectangle/chair dim-$6$ block and any fixed basis transported from it,
and more generally to any finite truncation once the remaining diagonal blocks are added.
\end{corollary}

\begin{proof}
Apply Proposition~\ref{prop:LaneA_triangular_assembly}; a fixed polynomial in polylogarithms is again polylogarithmic.
\end{proof}

\begin{remark}[Role of the explicit sample certificates]\label{rem:LaneA_sample_certificates}
The explicit CP-even scalar, CP-odd scalar, small-loop, traceless-tensor, and concrete dim-$6$ loop blocks now serve two referee-facing purposes.
First, they audit the one-step multiplier and UV-matching machinery in sectors where everything can be written out concretely.
Second, they provide normalization cross-checks for the canonical block protocol used later.
They are \emph{not} a second breadth theorem parallel to the canonical finite-truncation result below;
the load-bearing widening step is the passage from arbitrary fixed symmetry blocks to arbitrary finite engineering caps.
\end{remark}

\paragraph{Canonical flow-lifted blocks and arbitrary finite truncations.}
The explicit scalar/tensor/loop certificates above are useful audited examples, but the Lane~A inverse problem is not inherently restricted to those examples.
For the graded basis of local operators from Definition~\ref{def:local_operator_spaces}, one chooses a \emph{canonical flowed lift} block-by-block:
replace each curvature factor $F_{\mu\nu}$ by the flowed curvature $G_{\mu\nu}(t)$, each covariant derivative by its flowed counterpart,
project to the same $O(4)\times CP$ symmetry type, and impose the same scale-independent normalization moments that fix the graded triangular renormalization matrix of Definition~\ref{def:graded_Z}.
The point is not merely that the resulting matrix ``looks like the identity'' in some coordinates.
Appendix~\ref{app:normalization_crosscheck} proves a basis-free theorem on the exact-dimension quotient block itself:
the small-flow substitution acts as the identity on that quotient at tree level, and Theorem~\ref{thm:LaneA_block_coefficient_theorem} converts that quotient statement into an explicit coefficient-extraction theorem on every arbitrary exact-dimension block.
The proposition below upgrades that theorem-level tree statement to the quantitative $I+O(u(\mu_t))$ UV theorem used by the generic finite-dimensional diagonal-block criterion.

\begin{proposition}[Basis-free exact-dimension UV theorem for an arbitrary fixed symmetry block]\label{prop:canonical_block_uv_identity}
Fix a symmetry block $\mathfrak s$ in the graded basis of Definition~\ref{def:local_operator_spaces}, with engineering dimension $\Delta(\mathfrak s)$,
finite block size $m_{\mathfrak s}$, and quotient block $\mathscr V_{\mathfrak s}$ equipped with the canonical block norm of Definition~\ref{def:LaneA_canonical_block_norm}.
Let
\[
\Sigma_{\mathfrak s}^{\mathrm{UV}}(t):\mathscr V_{\mathfrak s}\longrightarrow \mathscr V_{\mathfrak s}
\]
denote the exact-dimension diagonal matching operator from the canonical flowed quotient block to the coherently renormalized quotient block at the matching scale $\mu_t\simeq 1/\sqrt t$.
Then
\begin{equation}\label{eq:canonical_block_uv_identity}
\Sigma_{\mathfrak s}^{\mathrm{UV}}(t)=I_{\mathfrak s}+E_{\mathfrak s}(t),
\qquad
\|E_{\mathfrak s}(t)\|_{\mathrm{op,can}}\le C_{\mathfrak s,\mathrm{UV}}\,u(\mu_t)
\end{equation}
for all sufficiently small $t$.
In particular, after shrinking $t_0$ so that $C_{\mathfrak s,\mathrm{UV}}u(\mu_t)\le \tfrac12$ on $(0,t_0]$, the UV block is invertible and
\begin{equation}\label{eq:canonical_block_uv_inv}
\|\big(\Sigma_{\mathfrak s}^{\mathrm{UV}}(t)\big)^{-1}\|_{\mathrm{op,can}}\le 2.
\end{equation}
If one writes $\Sigma_{\mathfrak s}^{\mathrm{UV}}(t)$ in the fixed canonical basis
\[
\mathbf Y_{\mathfrak s,t}(f)
=
\bigl(Y^{(1)}_{\mathfrak s,t}(f),\dots,Y^{(m_{\mathfrak s})}_{\mathfrak s,t}(f)\bigr)^{\mathsf T},
\qquad
\mathbf O_{\mathfrak s}^{\mathrm{ren}}(f;\mu)
=
\bigl(O^{(1),\mathrm{ren}}_{\mathfrak s}(f;\mu),\dots,O^{(m_{\mathfrak s}),\mathrm{ren}}_{\mathfrak s}(f;\mu)\bigr)^{\mathsf T},
\]
then the small-flow-time expansion has the block form
\begin{equation}\label{eq:canonical_block_uv_form}
\mathbf Y_{\mathfrak s,t}(f)
=
C_{\mathfrak s,\mathrm{UV}}(t)\,\mathbf O_{\mathfrak s}^{\mathrm{ren}}(f;\mu_t)
+\sum_{\mathfrak s'\prec \mathfrak s} B_{\mathfrak s\mathfrak s'}(t)\,\mathbf O_{\mathfrak s'}^{\mathrm{ren}}(f;\mu_t)
+\mathbf R_{\mathfrak s,t}(f),
\end{equation}
where $C_{\mathfrak s,\mathrm{UV}}(t)$ is exactly the matrix of $\Sigma_{\mathfrak s}^{\mathrm{UV}}(t)$ in that basis.
\end{proposition}

\begin{proof}
Theorem~\ref{thm:LaneA_block_coefficient_theorem} already identifies the exact-dimension diagonal tree-level matching matrix on
$\mathscr V_{\mathfrak s}$ as the identity in the canonical coordinates, with representative-independence and cap-independence built into the quotient-block formulation.
Thus the exact-dimension diagonal matching operator has tree-level part $I_{\mathfrak s}$.

To upgrade from tree level to the physical AF matching scale, use the insertion decomposition at matching scale (Theorem~\ref{thm:ins_decomp}) together with
Proposition~\ref{prop:Eins_bound}: on a fixed finite block, every exact-dimension diagonal coefficient is analytic in the AF coupling and differs from its tree-level value by $O(u(\mu_t))$.
Because $m_{\mathfrak s}<\infty$, the canonical block norm is equivalent to any coefficient norm on that block, so these coefficientwise $O(u(\mu_t))$ bounds assemble into the operator bound in \eqref{eq:canonical_block_uv_identity}.
The inverse estimate \eqref{eq:canonical_block_uv_inv} is the Neumann-series bound for $I_{\mathfrak s}+E_{\mathfrak s}(t)$.
Finally, once the fixed canonical basis is chosen, the operator $\Sigma_{\mathfrak s}^{\mathrm{UV}}(t)$ is represented by the matrix $C_{\mathfrak s,\mathrm{UV}}(t)$, and the full SFTE takes the block form \eqref{eq:canonical_block_uv_form} with the lower-block terms and remainder supplied by Theorem~\ref{thm:ins_decomp}.
\end{proof}

\begin{corollary}[Canonical block matrix theorem]\label{cor:canonical_block_matrix_theorem}
Under the hypotheses of Proposition~\ref{prop:canonical_block_uv_identity},
the exact-dimension diagonal UV matrix in the canonical basis satisfies
\[
C_{\mathfrak s,\mathrm{UV}}(t)=I_{m_{\mathfrak s}}+E^{\mathrm{mat}}_{\mathfrak s}(t),
\qquad
\|E^{\mathrm{mat}}_{\mathfrak s}(t)\|_{\mathrm{op}}\le C_{\mathfrak s,\mathrm{UV}}\,u(\mu_t),
\]
and therefore obeys the coordinate version of \eqref{eq:canonical_block_uv_inv}.
In particular, the later finite-truncation inverse theorem may be read either intrinsically on the quotient blocks or, equivalently, in the chosen canonical matrices.
\end{corollary}

\begin{proof}
This is just Proposition~\ref{prop:canonical_block_uv_identity} written in the fixed canonical basis.
\end{proof}

\begin{remark}[Scope of the canonical UV identity]\label{rem:canonical_block_uv_scope}
Proposition~\ref{prop:canonical_block_uv_identity} concerns only the \emph{exact-dimension diagonal block}.
It does not assert that lower-dimension mixings vanish.
Those lower-dimension contributions are exactly the strictly lower blocks in \eqref{eq:canonical_block_uv_form};
they are handled later by the lower-triangular assembly theorem.
The point of the proposition is more specific:
once both sides are written in the same complete-contraction basis and the coherent normalization moments are fixed,
the map on the exact-dimension quotient is canonically the identity at tree level.
That is the bookkeeping fact which lets the later truncation theorem be stated without enumerating sectors one by one.
\end{remark}

\begin{theorem}[Canonical symmetry-block inverse control]\label{thm:canonical_block_polylog}
Fix any block $\mathfrak s$ in the graded local operator basis of Definition~\ref{def:local_operator_spaces} and let
$C_{\mathfrak s}(t,\mu)$ be the corresponding coefficient matrix from the canonical flowed lifts to the renormalized local basis at scale $\mu$.
Assume the AF-ball hypothesis along the matching trajectory.
Then for all sufficiently small $t$ the block is invertible and obeys
\begin{equation}\label{eq:canonical_block_polylog}
\|C_{\mathfrak s}(t,\mu)^{-1}\|_{\mathrm{op}}
\le
C_{\mathfrak s,\mathrm{inv}}
\Big(\frac{u(\mu)}{u(\mu_t)}\Big)^{p_{\mathfrak s}},
\qquad
p_{\mathfrak s}:=\frac{4C_{M,\mathfrak s}}{A_-},
\end{equation}
so the inverse growth is again at most polylogarithmic in $t^{-1}$.
\end{theorem}

\begin{proof}
Combine Proposition~\ref{prop:block_mult_bound} with Proposition~\ref{prop:canonical_block_uv_identity} and apply Theorem~\ref{thm:block_polylog} with $C_{\mathfrak s}^{(0)}=I_{m_{\mathfrak s}}$.
\end{proof}

\begin{theorem}[Basis-independent finite-truncation inverse control for the full sharp-local basis]\label{thm:finite_truncation_inverse_control}
Fix $\Delta_{\max}\ge 0$ and order the canonical symmetry blocks occurring in the graded basis of Definition~\ref{def:local_operator_spaces}
with engineering dimension $\le \Delta_{\max}$ by increasing engineering dimension and then by a fixed symmetry-type order.
Let
\[
C^{(\le \Delta_{\max})}(t,\mu)
\]
be the full coefficient matrix from the canonical flowed lifts of that truncation to the corresponding renormalized local basis.
Then, for all sufficiently small $t$:
\begin{enumerate}[label=(\roman*),leftmargin=2.3em]
\item $C^{(\le \Delta_{\max})}(t,\mu)$ is block lower triangular in the chosen grading order;
\item each diagonal block satisfies the polylogarithmic inverse bound of Theorem~\ref{thm:canonical_block_polylog};
\item each off-diagonal block is uniformly bounded on $(0,t_0]$ by the insertion integration and insertion-to-OS bounds (Proposition~\ref{prop:Eins_bound} and Theorem~\ref{thm:ins_to_OS}) together with finite-depth composition in the fixed truncation.
\end{enumerate}
Consequently there exist constants $C_{\Delta_{\max},\mathrm{inv}}<\infty$ and $p_{\Delta_{\max}}<\infty$ such that
\begin{equation}\label{eq:finite_truncation_polylog}
\big\|\big(C^{(\le \Delta_{\max})}(t,\mu)\big)^{-1}\big\|_{\mathrm{op}}
\le
C_{\Delta_{\max},\mathrm{inv}}
\Big(\frac{u(\mu)}{u(\mu_t)}\Big)^{p_{\Delta_{\max}}}
\qquad (0<t\le t_0).
\end{equation}
In particular, the Lane~A quantitative inverse-control hypothesis \eqref{eq:LaneA_quant_inverse_bound} holds for every finite engineering-dimension truncation of the sharp-local operator basis.
\end{theorem}

\begin{proof}
Power counting gives the block lower-triangular structure in the grading order.
Each diagonal block obeys canonical polylogarithmic inverse control by Theorem~\ref{thm:canonical_block_polylog}.
For the finitely many off-diagonal blocks, Proposition~\ref{prop:Eins_bound} and Theorem~\ref{thm:ins_to_OS} give a scale-uniform control on the corresponding insertion coefficients; because the truncation is finite, repeated composition across the finite RG depth only changes the constant.
Apply Proposition~\ref{prop:LaneA_triangular_assembly} to the resulting block lower-triangular system and absorb the finite polynomial in the diagonal polylog factors into a single exponent $p_{\Delta_{\max}}$ and prefactor $C_{\Delta_{\max},\mathrm{inv}}$.
\end{proof}

\begin{corollary}[Lane~A closure on every finite truncation]\label{cor:LaneA_all_finite_truncations}
For every $\Delta_{\max}$, the finite family of all canonical flowed lifts / renormalized local operators of engineering dimension $\le \Delta_{\max}$ carries the intrinsic capwise data used in Theorem~\ref{thm:LaneA_closure}:
\begin{enumerate}[label=(\roman*),leftmargin=2.3em]
\item the quantitative inverse-control package of Theorem~\ref{thm:finite_truncation_inverse_control};
\item the canonical SFTE / remainder-compatibility package of Corollary~\ref{cor:finite_truncation_SFTE}.
\end{enumerate}
Therefore Theorem~\ref{thm:LaneA_closure} applies canonically and without extra sectorwise hypotheses to every finite engineering-dimension truncation of the sharp-local algebra.
\end{corollary}

\begin{proof}
This is exactly the combination of Theorem~\ref{thm:finite_truncation_inverse_control}, Corollary~\ref{cor:finite_truncation_SFTE}, and the finite-cap extension statement of Theorem~\ref{thm:LaneA_closure}.
\end{proof}

\begin{remark}[Current scope after the finite-truncation upgrade]
The explicit scalar/tensor/loop blocks above remain useful audited examples and normalization anchors.
But the canonical-block results now close Packet~5 at the level actually needed downstream:
every fixed symmetry block in the graded basis of Definition~\ref{def:local_operator_spaces} is controlled by the basis-free associated-graded identity theorem and hence has a canonical UV normalization,
every such block therefore satisfies the polylogarithmic inverse bound of Theorem~\ref{thm:canonical_block_polylog},
and every finite engineering-dimension truncation inherits both quantitative inverse control and SFTE/remainder compatibility by Theorem~\ref{thm:finite_truncation_inverse_control} and Corollary~\ref{cor:finite_truncation_SFTE}.
Table~\ref{tab:LaneA_packet_interface_map} in Appendix~\ref{app:normalization_crosscheck} records the exact role boundary: the present section exports the unconditional finite-truncation input used later by Theorem~\ref{thm:finite_cap_flowed_sector_bridge} and Theorem~\ref{thm:LaneA_closure}, but it does not by itself define the capwise mixed state, the inductive-union sharp-local state, or bounded-base cyclicity.
Appendix~\ref{app:normalization_crosscheck} is therefore the referee-facing dictionary for the chosen canonical bases and normalization conventions, not a second closure proof.
\end{remark}

